\chapter{Các công trình nghiên cứu liên quan}
\label{Chapter2}

Chương này khảo sát các hướng nghiên cứu có liên quan trực tiếp đến bài toán của luận văn, bao gồm: phân loại ảnh y khoa đa phương thức, thích nghi mô hình ngôn ngữ-thị giác trên dữ liệu y khoa hạn chế, xử lý mất cân bằng lớp và phát hiện dữ liệu ngoài phân phối. Thay vì liệt kê riêng lẻ từng công trình, nội dung được tổ chức theo tiến trình phát triển của các phương pháp, trong đó mỗi hướng tiếp cận được trình bày thông qua hạn chế mà nó giải quyết và vấn đề còn tồn tại. Từ đó, chương này hình thành cơ sở cho các lựa chọn thiết kế được trình bày ở Chương~\ref{Chapter3}.

\section{Phân loại ảnh y khoa đa phương thức}
\label{sec:related_multimodal}

\subsection{Từ ghép nối đặc trưng đến tương tác liên phương thức}

Các nghiên cứu ban đầu thường kết hợp ảnh y khoa và dữ liệu lâm sàng bằng cách nối các véc-tơ đặc trưng rồi đưa biểu diễn hợp nhất vào một bộ phân loại. Chẳng hạn, X. Mei và cộng sự \cite{XMei} kết hợp đặc trưng ảnh từ CNN với thông tin lâm sàng thông qua một mạng MLP. Kết quả cho thấy văn bản có thể bổ sung những thông tin không thể hiện rõ trên hình ảnh. Tuy nhiên, phép nối đặc trưng chỉ tạo ra sự kết hợp ở mức toàn cục và chưa mô hình hóa được mối quan hệ giữa từng vùng ảnh với nội dung văn bản tương ứng.

Để khắc phục hạn chế này, TieNet \cite{TieNet} bổ sung cơ chế chú ý nhằm liên kết các từ quan trọng trong báo cáo với những vùng ảnh có liên quan. IRENE \cite{IRENE} tiếp tục mở rộng sự tương tác bằng cách biểu diễn hai phương thức dưới dạng chuỗi token và cho phép các token ảnh và văn bản trao đổi thông tin trong Transformer. Như vậy, quá trình phát triển của các mô hình đa phương thức đã chuyển từ ghép nối đặc trưng tĩnh sang học quan hệ liên phương thức ở mức chi tiết hơn.

Dù đạt hiệu quả tích cực, các kiến trúc được xây dựng và huấn luyện riêng cho từng tác vụ thường phụ thuộc vào số lượng lớn dữ liệu ảnh-văn bản được ghép cặp, đồng nghĩa với việc có chi phí huấn luyện cao. Hạn chế này thúc đẩy việc kế thừa các mô hình ngôn ngữ-thị giác đã được tiền huấn luyện trên dữ liệu quy mô lớn.

\subsection{Sự phát triển của các mô hình ngôn ngữ-thị giác y khoa}
\label{subsec:med_vlm_progress}

Các mô hình ngôn ngữ-thị giác y khoa (Medical Vision-Language Models, Med-VLMs) sử dụng các bộ mã hóa riêng cho ảnh và văn bản, sau đó căn chỉnh hai không gian biểu diễn bằng các mục tiêu tiền huấn luyện đa phương thức. Cách tiếp cận này cho phép tái sử dụng tri thức từ dữ liệu quy mô lớn và giảm sự phụ thuộc vào nhãn phân loại thủ công của từng tác vụ đích.

ConVIRT \cite{ConVIRT} là một trong những công trình đầu tiên áp dụng học tương phản để căn chỉnh biểu diễn toàn cục của ảnh y khoa và báo cáo văn bản. Phương pháp này chứng minh khả năng khai thác dữ liệu ảnh-báo cáo mà không yêu cầu nhãn phân loại chi tiết. Tuy nhiên, căn chỉnh ở mức toàn cục chưa phản ánh đầy đủ quan hệ giữa vùng tổn thương và từ ngữ mô tả tương ứng.

GLoRIA \cite{GLoRIA} giải quyết một phần hạn chế trên bằng cách kết hợp mục tiêu tương phản toàn cục và cục bộ, qua đó học quan hệ giữa các vùng ảnh và token văn bản. Mặc dù tạo ra biểu diễn chi tiết hơn, các phương pháp tương phản theo cặp vẫn có thể xem những mẫu khác nhau nhưng cùng biểu hiện bệnh lý là cặp âm tính, dẫn đến hiện tượng âm tính giả.

MedCLIP \cite{MedCLIP} giảm sự phụ thuộc vào các cặp ảnh-văn bản tương ứng tuyệt đối bằng cơ chế tách cặp và hàm mất mát Semantic Matching. Thay vì coi toàn bộ các cặp không trùng khớp là âm tính hoàn toàn, phương pháp khai thác mức độ tương đồng ngữ nghĩa giữa thông tin y khoa của các mẫu. Tuy nhiên, phạm vi tiền huấn luyện của MedCLIP chủ yếu tập trung vào ảnh X-quang ngực.

BioMedCLIP \cite{BiomedCLIP} mở rộng quy mô và phạm vi tiền huấn luyện bằng tập PMC-15M, sử dụng ViT-Base cho hình ảnh và PubMedBERT cho văn bản. Nhờ được huấn luyện trên nhiều loại dữ liệu y sinh, BioMedCLIP cung cấp nền tảng phù hợp hơn để chuyển giao sang các tác vụ chuyên biệt. Dù vậy, việc áp dụng trực tiếp mô hình này vào ảnh X-quang xương vẫn gặp sự khác biệt về loại hình ảnh, loại tổn thương và nguồn văn bản đầu vào. Đặc biệt, dữ liệu của luận văn sử dụng bệnh sử lâm sàng được thu nhận trước khi đọc ảnh thay vì báo cáo mô tả trực tiếp nội dung X-quang. Do đó, mô hình nền tảng vẫn cần được thích nghi với miền dữ liệu đích.

Bảng~\ref{tab:medvlm_progress} tóm tắt vai trò của các công trình chính trong tiến trình phát triển trên.

\begin{table}[H]
\centering
\resizebox{\textwidth}{!}{
\begin{tabular}{|l|p{4.1cm}|p{5.2cm}|}
\hline
\textbf{Phương pháp} & \textbf{Hạn chế được giải quyết} & \textbf{Hạn chế còn lại} \\ \hline
ConVIRT \cite{ConVIRT} &
Khai thác dữ liệu ảnh-báo cáo bằng căn chỉnh tương phản, giảm nhu cầu nhãn phân loại thủ công. &
Chủ yếu căn chỉnh ở mức toàn cục, chưa mô hình hóa chi tiết quan hệ vùng ảnh-token. \\ \hline
GLoRIA \cite{GLoRIA} &
Bổ sung căn chỉnh toàn cục và cục bộ giữa vùng ảnh và từ ngữ. &
Vẫn chịu ảnh hưởng của các cặp âm tính giả trong học tương phản. \\ \hline
MedCLIP \cite{MedCLIP} &
Sử dụng Semantic Matching để giảm phụ thuộc vào cặp ảnh-văn bản tuyệt đối và hạn chế âm tính giả. &
Phạm vi tiền huấn luyện chủ yếu tập trung vào ảnh X-quang ngực. \\ \hline
BioMedCLIP \cite{BiomedCLIP} &
Mở rộng quy mô và miền y sinh của dữ liệu tiền huấn luyện. &
Chưa được thích nghi riêng cho ảnh X-quang xương và bệnh sử lâm sàng trước khi đọc ảnh. \\ \hline
\end{tabular}
}
\caption{Tiến trình giải quyết hạn chế của các mô hình ngôn ngữ-thị giác y khoa có liên quan đến luận văn.}
\label{tab:medvlm_progress}
\end{table}

Song song với Med-VLMs, các mô hình ngôn ngữ lớn đa phương thức như Med-Flamingo \cite{moor_2023_medflamingo}, LLaVA-Med \cite{li_2023_llavamed} và Med-PaLM M \cite{MED-PALM-M} mở rộng khả năng hội thoại và suy luận trên nhiều tác vụ y khoa. Tuy nhiên, các mô hình này có quy mô lớn, chi phí huấn luyện và suy luận cao, đồng thời tồn tại nguy cơ sinh nội dung không chính xác. Vì vậy, đối với một tác vụ phân loại chuyên biệt và có tập nhãn xác định trước, việc sử dụng Med-VLM kết hợp chiến lược tinh chỉnh hiệu quả tham số phù hợp hơn với phạm vi của luận văn.

\section{Thích nghi mô hình nền tảng trên dữ liệu y khoa hạn chế}
\label{sec:peft_related}

Tinh chỉnh toàn phần cập nhật toàn bộ tham số của mô hình nền tảng và thường đạt khả năng thích nghi cao, nhưng yêu cầu nhiều bộ nhớ, chi phí tính toán lớn và có nguy cơ làm suy giảm tri thức tiền huấn luyện khi dữ liệu đích có quy mô nhỏ. Ngược lại, dò tuyến tính chỉ huấn luyện lớp phân loại cuối nên tiết kiệm tài nguyên, nhưng không đủ linh hoạt khi tồn tại chênh lệch miền lớn giữa dữ liệu tiền huấn luyện và dữ liệu đích.

Tinh chỉnh hiệu quả tham số (Parameter-Efficient Fine-Tuning, PEFT) giải quyết sự đánh đổi này bằng cách chỉ tối ưu một phần nhỏ tham số. Các hướng phổ biến gồm tinh chỉnh câu nhắc, bộ điều hợp và tái tham số hóa hạng thấp. Tinh chỉnh câu nhắc thay đổi đầu vào có thể học được nhưng nhạy với cách khởi tạo và bị giới hạn bởi không gian đầu vào; các công trình như PLIP \cite{PLIP} và BiomedCoOp \cite{BiomedCoOp} cho thấy hiệu quả của hướng này trong điều kiện ít mẫu. Adapter Tuning chèn các mô-đun nhỏ giữa các lớp của mô hình, có tính mô-đun cao nhưng làm tăng thành phần phải thực thi trong quá trình suy luận; CLIPath \cite{CLIPath} và FTB \cite{FTB} là các đại diện trong lĩnh vực y khoa.

\subsection{Tinh chỉnh bằng tái tham số hóa hạng thấp}
\label{subsec:lora}

LoRA (Low-Rank Adaptation) đóng băng ma trận trọng số đã được tiền huấn luyện và biểu diễn phần cập nhật bằng tích của hai ma trận hạng thấp. Với ma trận trọng số
$W \in \mathbb{R}^{d_{\mathrm{out}} \times d_{\mathrm{in}}}$, phép biến đổi sau khi bổ sung LoRA được viết như sau:

\begin{equation}
W' = W + \Delta W
   = W + \frac{\alpha}{r}BA,
\label{eq:lora_update}
\end{equation}

trong đó $A \in \mathbb{R}^{r \times d_{\mathrm{in}}}$,
$B \in \mathbb{R}^{d_{\mathrm{out}} \times r}$, $r$ là hạng của phép phân rã và $\alpha$ là hệ số tỷ lệ. Trong quá trình tinh chỉnh, chỉ các ma trận $A$ và $B$ được cập nhật. Vì $r$ nhỏ hơn nhiều so với kích thước của $W$, số lượng tham số cần huấn luyện giảm đáng kể. LCBB \cite{LCBB} cho thấy việc áp dụng LoRA cho mô hình nền tảng y khoa có thể duy trì hiệu năng gần với tinh chỉnh toàn phần trong điều kiện dữ liệu hạn chế.

\subsection{So sánh LoRA và QLoRA}
\label{subsec:lora_qlora}

QLoRA kết hợp cơ chế thích ứng hạng thấp với lượng tử hóa trọng số của mô hình nền. Trong khi LoRA lưu các trọng số đóng băng ở độ chính xác dấu phẩy động, QLoRA thường lưu phần lớn trọng số nền ở dạng 4-bit và chỉ huấn luyện các ma trận LoRA. Vì vậy, QLoRA chủ yếu giảm bộ nhớ dùng để tải mô hình nền, chứ không làm thay đổi bản chất của các tham số hạng thấp được tối ưu.

\begin{table}[H]
\centering
\begin{tabular}{|p{3.3cm}|p{5.0cm}|p{5.0cm}|}
\hline
\textbf{Tiêu chí} & \textbf{LoRA} & \textbf{QLoRA} \\ \hline
Trọng số mô hình nền &
Được đóng băng và lưu ở độ chính xác dấu phẩy động. &
Được đóng băng và lượng tử hóa, thường xuống 4-bit. \\ \hline
Tham số được huấn luyện &
Các ma trận thích ứng hạng thấp. &
Các ma trận thích ứng hạng thấp. \\ \hline
Mức sử dụng bộ nhớ &
Thấp hơn tinh chỉnh toàn phần nhưng cao hơn QLoRA. &
Thấp hơn LoRA do trọng số nền được lượng tử hóa. \\ \hline
Độ chính xác số học &
Không có sai số lượng tử hóa ở trọng số nền. &
Có thể chịu ảnh hưởng của sai số lượng tử hóa. \\ \hline
Độ phức tạp triển khai &
Đơn giản và ổn định hơn khi áp dụng đồng thời lên bộ mã hóa ảnh và văn bản. &
Phụ thuộc vào thư viện, phần cứng và toán tử hỗ trợ lượng tử hóa. \\ \hline
Trường hợp phù hợp &
Mô hình có quy mô vừa và phần cứng đủ để lưu trọng số nền. &
Mô hình rất lớn hoặc môi trường bị giới hạn nghiêm ngặt về bộ nhớ. \\ \hline
\end{tabular}
\caption{So sánh LoRA và QLoRA trong tinh chỉnh mô hình nền tảng.}
\label{tab:lora_qlora}
\end{table}

Trong luận văn, LoRA được lựa chọn thay cho QLoRA. BioMedCLIP sử dụng các bộ mã hóa ViT-Base và PubMedBERT, có quy mô nhỏ hơn đáng kể so với các mô hình ngôn ngữ lớn mà QLoRA thường hướng đến. Mô hình có thể được tải và tinh chỉnh bằng LoRA trong giới hạn bộ nhớ của phần cứng được sử dụng, nên lợi ích tiết kiệm bộ nhớ từ lượng tử hóa không đủ lớn để bù cho độ phức tạp bổ sung. Bên cạnh đó, việc giữ trọng số nền ở độ chính xác dấu phẩy động giúp tránh nguy cơ sai số lượng tử hóa ảnh hưởng đến các đặc trưng nhỏ và tinh tế trên ảnh X-quang xương. Do đó, LoRA tạo ra sự cân bằng phù hợp hơn giữa hiệu quả tham số, độ chính xác số học và độ ổn định triển khai.

\section{Phân loại trong điều kiện mất cân bằng dữ liệu}
\label{sec:wce_related}

Các bộ dữ liệu y khoa thường có sự chênh lệch lớn về số lượng mẫu giữa các lớp. Nếu sử dụng Cross-Entropy thông thường, các lớp phổ biến xuất hiện nhiều lần hơn trong hàm mục tiêu và có thể chi phối quá trình tối ưu. Hệ quả là mô hình đạt độ chính xác tổng thể cao nhưng nhận diện kém các lớp hiếm.

Weighted Cross-Entropy (WCE) khắc phục vấn đề này bằng cách gán trọng số lớn hơn cho các lớp có ít mẫu. Với bài toán phân loại $K$ lớp và $N$ mẫu, hàm mất mát được xác định như sau:

\begin{equation}
\mathcal{L}_{\mathrm{WCE}}
= -\frac{1}{N}\sum_{i=1}^{N}
\omega_{y_i}\log p_{i,y_i},
\label{eq:wce}
\end{equation}

trong đó $p_{i,y_i}$ là xác suất mà mô hình gán cho lớp đúng $y_i$ của mẫu thứ $i$, còn $\omega_{y_i}$ là trọng số của lớp tương ứng. Một cách xác định trọng số thường dùng là nghịch đảo tần suất đã chuẩn hóa:

\begin{equation}
\omega_k = \frac{N}{K n_k},
\label{eq:wce_weight}
\end{equation}

trong đó $n_k$ là số mẫu thuộc lớp $k$. Cách gán trọng số cụ thể cần được xác định trên tập huấn luyện và giữ cố định khi đánh giá để tránh rò rỉ thông tin từ tập kiểm thử.

So với Focal Loss, WCE tác động trực tiếp ở cấp lớp và có cách diễn giải đơn giản hơn: lớp hiếm nhận trọng số lớn hơn trong quá trình tối ưu. Focal Loss tập trung vào độ khó của từng mẫu và có thể hữu ích khi tồn tại nhiều mẫu dễ, nhưng yêu cầu thêm siêu tham số và có thể làm quá trình so sánh khó kiểm soát hơn. Vì mục tiêu chính của luận văn là giảm ảnh hưởng của mất cân bằng lớp, WCE được lựa chọn làm hàm mất mát phân loại trong giai đoạn huấn luyện mô-đun dung hợp và bộ phân loại.

\section{Phát hiện dữ liệu ngoài phân phối}
\label{sec:ood_related}

Phát hiện dữ liệu ngoài phân phối (Out-of-Distribution Detection, OOD Detection) nhằm nhận diện những mẫu không phù hợp với phân phối dữ liệu đã dùng để huấn luyện. Trong chẩn đoán hình ảnh y khoa, cơ chế này có ý nghĩa quan trọng vì mô hình có thể gặp bệnh lý hiếm, vùng giải phẫu không thuộc tập nhãn, ảnh có chất lượng bất thường hoặc các trường hợp chưa từng xuất hiện trong dữ liệu huấn luyện. Nếu không có cơ chế OOD, mô hình phân loại đóng vẫn buộc phải gán mẫu vào một lớp đã biết, kể cả khi dự đoán không đáng tin cậy.

\subsection{Phương pháp hậu xử lý dựa trên đầu ra}

Maximum Softmax Probability (MSP) sử dụng xác suất Softmax lớn nhất làm điểm tin cậy. Mẫu có xác suất cực đại thấp được xem là ứng viên OOD. Phương pháp này đơn giản và không cần thay đổi mô hình, nhưng mạng nơ-ron sâu có thể đưa ra xác suất cao ngay cả trên dữ liệu khác phân phối.

Energy-based Score sử dụng trực tiếp vector logit thay vì xác suất đã chuẩn hóa. Với logit $f_k(x)$ và hệ số nhiệt độ $T$, mức năng lượng được tính bởi:

\begin{equation}
E(x) = -T\log\sum_{k=1}^{K}
\exp\left(\frac{f_k(x)}{T}\right).
\label{eq:energy_score}
\end{equation}

Các mẫu trong phân phối thường có năng lượng thấp hơn, trong khi mẫu OOD có xu hướng nhận năng lượng cao hơn. Tuy nhiên, cả MSP và Energy Score đều phụ thuộc mạnh vào đặc tính của đầu ra bộ phân loại và mức hiệu chỉnh xác suất của mô hình.

\subsection{Phương pháp dựa trên không gian đặc trưng}

Nhóm phương pháp này sử dụng biểu diễn $z=h(x)$ được trích xuất từ mô hình và đánh giá mức độ phù hợp của biểu diễn đó với phân bố đặc trưng của dữ liệu huấn luyện. Khoảng cách Euclidean chỉ xét độ lớn sai khác và ngầm giả định phân bố có phương sai đồng đều theo mọi chiều. Khoảng cách Cosine chỉ xét hướng của hai véc-tơ và bỏ qua cấu trúc phương sai của dữ liệu. Trong khi đó, khoảng cách Mahalanobis sử dụng ma trận hiệp phương sai để điều chỉnh đóng góp của từng chiều và mối tương quan giữa các chiều đặc trưng.

Để phù hợp với cơ chế hậu xử lý được sử dụng trong luận văn, mỗi lớp $k$ được biểu diễn bởi trung bình đặc trưng:

\begin{equation}
\boldsymbol{\mu}_k
= \frac{1}{N_k}
\sum_{i:y_i=k}\boldsymbol{z}_i,
\label{eq:maha_mean}
\end{equation}

trong đó $N_k$ là số mẫu huấn luyện thuộc lớp $k$. Thay vì ước lượng một ma trận hiệp phương sai riêng cho từng lớp, một ma trận hiệp phương sai gộp dùng chung được tính trên toàn bộ dữ liệu huấn luyện:

\begin{equation}
\boldsymbol{\Sigma}
= \frac{1}{N-K}
\sum_{k=1}^{K}
\sum_{i:y_i=k}
(\boldsymbol{z}_i-\boldsymbol{\mu}_k)
(\boldsymbol{z}_i-\boldsymbol{\mu}_k)^{\top}
+ \lambda\boldsymbol{I},
\label{eq:maha_cov}
\end{equation}

trong đó $\lambda\boldsymbol{I}$ là thành phần điều chuẩn giúp ma trận ổn định và khả nghịch. Khoảng cách Mahalanobis bình phương giữa mẫu kiểm thử $\boldsymbol{z}$ và lớp $k$ được xác định bởi:

\begin{equation}
d_k(\boldsymbol{z})
= (\boldsymbol{z}-\boldsymbol{\mu}_k)^{\top}
\boldsymbol{\Sigma}^{-1}
(\boldsymbol{z}-\boldsymbol{\mu}_k).
\label{eq:maha_class_distance}
\end{equation}

Điểm OOD của mẫu là khoảng cách nhỏ nhất đến các phân bố lớp đã biết:

\begin{equation}
s_{\mathrm{Maha}}(\boldsymbol{z})
= \min_{k\in\{1,\ldots,K\}} d_k(\boldsymbol{z}).
\label{eq:maha_score}
\end{equation}

Mẫu được xác định là OOD khi
$s_{\mathrm{Maha}}(\boldsymbol{z})>\tau$, với $\tau$ là ngưỡng được lựa chọn trên tập xác thực. Công thức trên sử dụng khoảng cách bình phương và ma trận hiệp phương sai gộp, phù hợp với cơ chế hậu xử lý của pipeline; vì vậy không sử dụng căn bậc hai hoặc ma trận hiệp phương sai riêng cho từng lớp.

\section{Khoảng trống nghiên cứu và định hướng của đề tài}
\label{sec:chapter2_gap}

Qua các công trình đã khảo sát, có thể rút ra bốn khoảng trống chính đối với bài toán phân loại ảnh X-quang xương đa phương thức.

Thứ nhất, các phương pháp dung hợp ban đầu đã chuyển từ phép nối đặc trưng sang tương tác vùng ảnh-token, nhưng phần lớn được phát triển trên các miền dữ liệu khác và thường sử dụng báo cáo đọc ảnh. Việc khai thác bệnh sử lâm sàng được thu nhận trước khi đọc ảnh vẫn chưa được nghiên cứu đầy đủ, trong khi loại văn bản này có ít thông tin mô tả trực tiếp tổn thương hơn và tạo ra chênh lệch ngữ nghĩa lớn với dữ liệu tiền huấn luyện.

Thứ hai, BioMedCLIP cung cấp biểu diễn y sinh tổng quát nhưng chưa được thích nghi riêng cho ảnh X-quang xương và bệnh sử lâm sàng. Tinh chỉnh toàn phần có chi phí cao và dễ quá khớp trên tập dữ liệu nhỏ; QLoRA tiết kiệm thêm bộ nhớ nhưng tạo ra độ phức tạp và sai số lượng tử hóa không cần thiết đối với quy mô của BioMedCLIP. Vì vậy, LoRA được lựa chọn để thích nghi các bộ mã hóa với miền dữ liệu đích.

Thứ ba, phân bố lớp của dữ liệu thực tế có sự mất cân bằng đáng kể. Do pipeline không sử dụng Prototypical Loss, WCE được dùng làm hàm mất mát phân loại nhằm tăng đóng góp của các lớp hiếm mà không bổ sung một mục tiêu metric-learning riêng.

Cuối cùng, bộ phân loại đóng không thể tự nhận biết một mẫu nằm ngoài các lớp đã học. Luận văn vì vậy áp dụng Mahalanobis theo cơ chế hậu xử lý trên không gian đặc trưng, sử dụng trung bình lớp và ma trận hiệp phương sai gộp được ước lượng từ tập huấn luyện. Cơ chế này tách biệt khỏi quá trình tối ưu phân loại và cho phép xác định ngưỡng OOD trên tập xác thực.

Từ các phân tích trên, pipeline của luận văn được định hướng theo hai giai đoạn. Ở giai đoạn thứ nhất, các bộ mã hóa của BioMedCLIP được thích nghi bằng LoRA và mục tiêu Semantic Matching nhằm giảm chênh lệch miền giữa ảnh X-quang xương và bệnh sử lâm sàng. Ở giai đoạn thứ hai, biểu diễn ảnh và văn bản được dung hợp bằng các mô-đun tương tác liên phương thức, sau đó bộ phân loại được tối ưu bằng WCE. Sau huấn luyện, Mahalanobis được áp dụng theo cơ chế hậu xử lý để phát hiện dữ liệu ngoài phân phối. Kiến trúc và quy trình tối ưu cụ thể được trình bày trong Chương~\ref{Chapter3}.
